\section{交换环的一般性质}\label{sec:galois-1-2}
除非另有说明，本节中的环都假定为交换环.

\subsection{零环}\label{sec:galois-1-2-1}
对于一个（带单位的）环 $(A,+,\cdot)$，分配律蕴含，对任意 $a$ 都有
$a\cdot0=a\cdot(0+0)=a\cdot0+a\cdot0$，从而 $a\cdot0=0$.
因此，如果 $0=1$（我们并未排除这种情形），那么 $A=\{0\}$.
当然，可以排除这种病态情形；不过，在讨论商时保留这种情形会更方便.

\subsection{子环}\label{sec:galois-1-2-2}
环 $A$ 的子集 $B$ 称为\newterm{galois-subring}{子环}{subring}，如果它非空、对减法和乘法封闭，并且包含 $1$.

\subsection{零因子、正则元与整环}\label{sec:galois-1-2-3}
环中的非零元素 $a$ 称为\newterm{galois-zero-divisor}{零因子}{zero divisor}，如果存在非零元素 $a'$ 使得 $aa'=0$.
非零且不是零因子的元素 $a$ 称为\newterm{galois-regular-element}{正则元}{regular element}.
整环是没有零因子的环，也就是说，它满足 $ab=0\Rightarrow a=0\text{ 或 }b=0$.
在整环中，即使 $a$ 不可逆，也可以在等式中作消去：
\begin{align*}
a\ne0\text{ 且 }ab=ac\quad\Longrightarrow\quad b=c.
\end{align*}
\begin{example*}
显然，整环的子环仍是整环. 另一方面，任何域显然都是整环.
因此，代数整数环总是整环.
\end{example*}
\begin{example*}
考虑 $\mathbb C^2$ 中由方程 $X_1X_2=0$ 定义的代数子集 $V$.
它是两条坐标轴 $X_1=0$ 和 $X_2=0$ 的并.
多项式函数 $X_1:\mathbb C^2\longrightarrow\mathbb C$，$(z_1,z_2)\longmapsto z_1$，在 $V$ 上诱导出一个显然非零的多项式函数 $x_1$.
同样，$X_2$ 在 $V$ 上诱导出一个非零函数 $x_2$.
但由 $V$ 的定义，函数 $x_1x_2$ 在 $V$ 上处处为零.
因此，$x_1$ 和 $x_2$ 都是环 $\mathcal O(V)$ 中的零因子，所以这个环不是整环.
\end{example*}
\begin{exercise*}
对于哪些 $N>0$，环 $\mathbb Z/N\mathbb Z$ 是整环？
\end{exercise*}
\begin{exercise*}
设 $K$ 为域. 证明：有限维的交换 $K$-代数 $A$ 如果是整环，那么它是域.
将 $A$ 中乘以 $a\ne0$ 的映射看作 $A$ 的一个 $K$-线性自同态.
\end{exercise*}

\subsection{幂零元与既约环}\label{sec:galois-1-2-4}
元素 $x\in A$ 称为\newterm{galois-nilpotent-element}{幂零元}{nilpotent element}，如果存在整数 $k\in\mathbb N$ 使得 $x^k=0$.
特别地，非零幂零元是零因子.
使 $x^k=0$ 成立的最小整数 $k$ 称为 $x$ 的\newterm{galois-nilpotency-index}{幂零指数}{nilpotency index}.
一个环称为既约环，如果它没有非零幂零元.
因此，对环而言，有
\begin{align*}
\text{整环}\quad\Longrightarrow\quad\text{既约环}.
\end{align*}
\begin{example*}
考察 $\mathbb Z/N\mathbb Z$ 的情形.
如果 $N=p^\nu$，其中 $p$ 为素数且 $\nu>0$ 为整数，那么由定义，$p$ 在 $\mathbb Z/N\mathbb Z$ 中是幂零指数为 $\nu$ 的幂零元.
于是可以看到：
\begin{itemize}
\item 若 $\nu=1$，则 $\mathbb Z/p^\nu\mathbb Z$ 是域（从而是整环，也是既约环）.
\item 若 $\nu>1$，则 $\mathbb Z/p^\nu\mathbb Z$ 不是既约环，它的幂零元全体为 $p\mathbb Z/p^\nu\mathbb Z$.
\end{itemize}
更一般地，作分解 $N=\prod_p p^{\nu_p(N)}$，并利用下面将回顾的中国剩余定理，可见 $\mathbb Z/N\mathbb Z$ 是既约环，当且仅当 $N$ 不含平方因子，也就是说，对所有 $p$ 都有 $\nu_p(N)=1$.
\footnote{译注：原文写作``对所有 $p$ 都有 $\nu_p(N)=1$''. 此处应将 $p$ 的范围理解为整除 $N$ 的素数；不整除 $N$ 的素数满足 $\nu_p(N)=0$.}
\end{example*}

\subsection{环的直积}\label{sec:galois-1-2-5}
设 $A$ 和 $A'$ 为两个环. 在笛卡尔积 $A\times A'$ 上，规定
\begin{align*}
\begin{aligned}
(a,a')+(b,b')&:=(a+b,a'+b'),\\
(a,a')\cdot(b,b')&=(ab,a'b'),
\end{aligned}
\end{align*}
便赋予它一个环结构，称为\newterm{galois-product-ring}{直积环}{product ring}.
加法单位元为 $(0,0)$，乘法单位元为 $(1,1)$.
若两个环都不是零环，则直积 $A\times A'$ 不是整环，因为 $(1,0)\cdot(0,1)=(0,0)$.
\begin{example*}
中国剩余定理告诉我们：当 $\gcd(n,m)=1$ 时，乘积映射
\begin{align*}
\begin{aligned}
\mathbb Z/nm\mathbb Z&\longrightarrow\mathbb Z/n\mathbb Z\times\mathbb Z/m\mathbb Z,\\
a\pmod{nm}&\longmapsto(a\pmod n,a\pmod m)
\end{aligned}
\end{align*}
将 $\mathbb Z/nm\mathbb Z$ 与 $\mathbb Z/n\mathbb Z$ 和 $\mathbb Z/m\mathbb Z$ 的直积等同起来.
\end{example*}

\subsection{幂等元}\label{sec:galois-1-2-6}
在环 $A$ 中，若元素 $e$ 满足 $e^2=e$，就称它为\newterm{galois-idempotent}{幂等元}{idempotent}.
此时，子集 $eAe$ 继承一个环结构，其加法和乘法由 $A$ 中的运算诱导，乘法单位元为 $e$.
\begin{remark*}
$eAe$ 不是 $A$ 的子环，因为它的乘法单位元与 $A$ 的不同（除非 $e=1$）.
\end{remark*}
当 $A$ 交换时（更一般地，当 $e$ 是中心元，即与 $A$ 的所有元素交换时），我们直接有 $eAe=Ae$.
元素 $1-e$ 也是 $A$ 的一个幂等元，并且有理想的直和分解
$A=Ae\oplus A(1-e)$（注意，乘以 $e$ 的映射是一个投影算子，与线性代数中遇到的相同）.
这一分解将 $A$ 与环 $Ae$ 和 $A(1-e)$ 的直积等同起来.
更确切地，映射
\begin{align*}
A\longrightarrow Ae\times A(1-e),\qquad a\longmapsto(ae,a(1-e))
\end{align*}
是一个环同构，这一概念将在下面回顾. 它的逆映射为 $(x,y)\longmapsto x+y$.
\begin{example*}
设整数 $n$ 和 $m$ 互素. 取 $u,v\in\mathbb Z$ 使得 $un+vm=1$.
将这个等式乘以 $un$，便得到 $(un)^2\equiv un\pmod{nm}$.
因此，$un$ 在 $\mathbb Z/nm\mathbb Z$ 中的像 $e$ 是一个幂等元.
此外，有
$(\mathbb Z/nm\mathbb Z)e=n\mathbb Z/nm\mathbb Z$，它与 $\mathbb Z/m\mathbb Z$ 同构（同构的含义见下文）；同样，
$(\mathbb Z/nm\mathbb Z)(1-e)=m\mathbb Z/nm\mathbb Z$ 与 $\mathbb Z/n\mathbb Z$ 同构.
这样便再次得到中国剩余定理
$\mathbb Z/nm\mathbb Z\simeq\mathbb Z/n\mathbb Z\times\mathbb Z/m\mathbb Z$.
\end{example*}
\begin{remark*}[{几何解释}]
我们将在后面看到，代数子集 $V$ 连通，当且仅当其函数代数 $\mathcal O(V)$ 的幂等元只有 $0$ 和 $1$.
更一般地，代数子集的连通分支与其代数 $\mathcal O(V)$ 的本原幂等元一一对应.
若一个幂等元不能进一步分解为两个非零幂等元之和，就称它为\newterm{galois-primitive-idempotent}{本原幂等元}{primitive idempotent}.
\end{remark*}

\subsection{（同态）态射}\label{sec:galois-1-2-7}
环态射是一个保持环结构的映射 $\varphi:A\longrightarrow A'$，这里保持环结构是指：
\begin{itemize}
\item 对任意 $a,a'\in A$，有 $\varphi(a+a')=\varphi(a)+\varphi(a')$ 且 $\varphi(aa')=\varphi(a)\varphi(a')$.
\item $\varphi(1)=1$.
\end{itemize}
注意，条件 $\varphi(1)=1$ 并非无关紧要.
例如，映射 $\varphi:A\longrightarrow A\times A$，$a\longmapsto(a,0)$，满足第一条性质，却将 $1$ 映到 $(1,0)$：它不是环态射.
用幂等元的语言，并沿用上一分目的记号，$Ae$ 到 $A$ 的包含映射不是环态射.
\begin{remark*}
环态射 $A\longrightarrow A'$ 的像是 $A'$ 的子环.
\end{remark*}
\begin{remark*}
对任意环 $A$，存在唯一的环态射 $\mathbb Z\longrightarrow A$.
它将 $n\in\mathbb Z$ 映到 $n\times1$（其中 $n\times-$ 表示在 Abel 群 $A$ 中乘以 $n$）.
\end{remark*}
与通常的定义一样，环同构 $\varphi:A\longrightarrow A'$ 是一个存在双侧逆的态射；也就是说，存在态射 $\psi:A'\longrightarrow A$，使得 $\varphi\circ\psi=\mathrm{id}_{A'}$ 且 $\psi\circ\varphi=\mathrm{id}_A$.
\begin{lemma*}
一个态射是同构，当且仅当它作为映射是双射.
\end{lemma*}
\begin{proof}
同构显然是双射. 反过来，设 $\varphi$ 为双射；只需验证其逆双射 $\varphi^{-1}$ 是环态射. 这是直接的验证.
\end{proof}
\begin{exercise*}
验证直积 $A\times A'$ 到其任一因子 $A$ 或 $A'$ 的投影都是环态射.
若 $\psi:B\longrightarrow A$ 和 $\psi':B\longrightarrow A'$ 是两个环态射，验证
\begin{align*}
(\psi,\psi'):B\longrightarrow A\times A',\qquad b\in(\psi(b),\psi'(b))
\end{align*}
是环态射.\footnote{译注：原文在此写作 $b\in(\psi(b),\psi'(b))$；按上下文，此处表示映射规则 $b\mapsto(\psi(b),\psi'(b))$.}
证明每个态射 $\Psi:B\longrightarrow A\times A'$ 都具有这种形式.
\end{exercise*}
\begin{definition*}
设 $A$ 为交换环. 一个 $A$-代数是由环 $B$ 和环态射 $\psi:A\longrightarrow B$ 组成的二元组 $(B,\psi)$.
从 $(B,\psi)$ 到 $(B',\psi')$ 的 $A$-代数态射是一个满足 $\varphi\circ\psi=\psi'$ 的环态射 $\varphi:B\longrightarrow B'$.
\end{definition*}
\begin{remark*}
这个定义推广了域上代数的概念. 我们有时会滥用记号，省略其中的 $\psi$.
例如，直接说``设 $B$ 为一个 $\mathbb C$-代数''，而不说``设 $(B,\psi)$ 为一个 $\mathbb C$-代数''.
\end{remark*}
\begin{example*}
每个环都以唯一的方式成为一个 $\mathbb Z$-代数.
\end{example*}


\subsection{理想}\label{sec:galois-1-2-8}
前面已经回顾了环 $A$ 的理想的定义. 特别地，$A$ 是自身的一个理想. 若一个理想不同于 $A$，则称它为真理想. 此外，$\{0\}$ 也是理想，称为零理想. 理想的主要来源是下面的引理.
\begin{lemma*}
环同态 $\varphi\colon A\longrightarrow A'$ 的核 $\operatorname{Ker}(\varphi):=\varphi^{-1}(\{0\})$ 是 $A$ 的一个理想.
\end{lemma*}
\begin{proof}
由群论可知 $\operatorname{Ker}(\varphi)$ 是 $A$ 的加法子群. 因而只须验证它在乘以 $A$ 的元素后仍封闭. 若 $a\in\operatorname{Ker}(\varphi)$ 且 $a'\in A$，则 $\varphi(a'\cdot a)=\varphi(a')\cdot0=0$，故 $a'\cdot a\in\operatorname{Ker}(\varphi)$.
\end{proof}
反过来，我们稍后会看到，$A$ 的每个理想都是一个以 $A$ 为源的环态射的核，而且可以取满态射.
\begin{exercise*}
证明更一般地，$A'$ 的理想 $I'$ 的原像 $\varphi^{-1}(I')$ 是 $A$ 的理想. 举反例说明，$A$ 的理想 $I$ 的像 $\varphi(I)$ 未必是 $A'$ 的理想. 不过，证明若 $\varphi$ 是满射，则理想的像 $\varphi(I)$ 是理想.
\end{exercise*}
\begin{example*}[{幂零根}]
$A$ 中幂零元组成的集合 $\mathcal N(A)$ 是一个理想，称为 $A$ 的\newterm{galois-nilradical}{幂零根}{nilradical}. 由于 $A$ 交换，$\mathcal N(A)$ 在乘以 $A$ 的元素后仍封闭是显然的；它对加法的封闭性可由二项式公式
\begin{align*}
(x+y)^n=\sum_k\binom nk x^ky^{n-k}
\end{align*}
看出：若 $n$ 大于 $x$ 与 $y$ 的幂零指数之和，则 $(x+y)^n=0$.
\end{example*}
\begin{exercise*}[{理想的根}]
设 $I$ 是环 $A$ 的一个理想. 令
\begin{align*}
\sqrt I:=\{x\in A,\ \exists k\in\mathbb N^*,\ x^k\in I\}.
\end{align*}
证明 $\sqrt I$ 是包含 $I$ 的理想（可以注意到 $\sqrt{\{0\}}=\mathcal N(A)$，并尝试改用前面的论证）. 当 $A=\mathbb Z$ 且 $I=N\mathbb Z$ 时，求 $\sqrt I$.
\end{exercise*}
由于两个理想的交仍是理想，可以谈论包含 $A$ 的子集 $E$ 的最小理想（相对于包含关系而言）：它是所有包含 $E$ 的理想之交. 称之为 $E$ 生成的理想. 具体地，它由所有具有如下形式的 $x\in A$ 组成：
\begin{align*}
x=a_1e_1+\cdots+a_re_r,
\end{align*}
其中 $r\in\mathbb N^*$，$e_i\in E$，$a_i\in A$. 当 $E=\{x_1,\ldots,x_n\}$ 时，也将这个理想记为 $(x_1,\ldots,x_n)$ 或 $Ax_1+\cdots+Ax_n$.
\begin{definition*}
交换环 $A$ 的理想 $I$ 称为：
\begin{itemize}
\item 有限生成的，若它由 $A$ 的有限个元素生成.
\item 主理想，若它由单个元素生成（也可说是单生成的）.
\item 极大理想，若它在 $A$ 的真理想（即不同于 $A$ 的理想）中相对于包含关系是极大的.
\item 素理想，若它是真理想，且 $\forall x,y\in A,\ xy\in I\Rightarrow(x\in I\text{ 或 }y\in I)$.
\item 根理想，若它是真理想，且 $\forall x\in A,\ (\exists k\in\mathbb N^*,x^k\in I)\Rightarrow x\in I$.
\end{itemize}
\end{definition*}
显然，``主''蕴含``有限生成''. 还有下面的蕴含关系.
\begin{lemma*}
对理想 $I$，有 $I\text{ 极大}\Rightarrow I\text{ 为素理想}\Rightarrow I\text{ 为根理想}$.
\end{lemma*}
\begin{proof}
设 $I$ 极大，取 $x,y\in A$ 使得 $xy\in I$，并考虑理想 $(x)+I$. 若这个理想是真理想，由 $I$ 的极大性，它等于 $I$，于是 $x\in I$. 若它不是真理想，则 $(x)+I=A$，所以可写成 $1=ax+i$，其中 $i\in I$，$a\in A$；于是 $y=axy+iy$，这表明 $y\in I$. 因而 $I$ 是素理想. 另一个蕴含是直接的.
\end{proof}
\begin{example*}
在 $\mathbb Z$ 中，每个理想都是主理想，所以具有 $n\mathbb Z$ 的形式，其中 $n\geqslant0$ 唯一确定. 当 $n\ne1$ 时，这个理想是真理想. 此时，它是素理想当且仅当 $n$ 是素数，此时它也是极大理想. 此外，它是根理想当且仅当 $n$ 无平方因子（习题）.\footnote{译注：原文允许 $n=0$，但 $(0)$ 也是 $\mathbb Z$ 的素理想；因此关于素理想的上述``当且仅当''在 $n=0$ 时有例外. 原文断言照录.}
\end{example*}
\begin{example*}
环 $A$ 是整环，当且仅当它的零理想 $I=\{0\}$ 是素理想.
\end{example*}
\begin{example*}
在环 $A=\mathbb C[X,Y]$ 中，理想 $(X)$ 是素理想，但不是极大理想，因为它包含于 $(X,Y)$. 后者则是极大理想. 事实上，对任意多项式 $f=f(X,Y)$，都有 $f\in f(0,0)+(X,Y)$，从而 $f(0,0)\in(f,X,Y)$. 所以，若 $f\notin(X,Y)$，则数 $f(0,0)$（视为零次多项式）非零，因而在 $\mathbb C[X,Y]$ 中可逆；理想 $(f,X,Y)$ 包含一个可逆元，所以等于 $\mathbb C[X,Y]$. 因此 $(X,Y)$ 不包含于任何真理想.\footnote{译注：原文末句如此；按上下文，意指它不严格包含于任何真理想.}
\end{example*}
\begin{remark*}
在环 $\mathbb C[X_1,\ldots,X_n]$ 中，有一条相互嵌套的素理想``链''
\begin{align*}
(0)\subset(X_1)\subset(X_1,X_2)\subset\cdots\subset(X_1,\ldots,X_n).
\end{align*}
这条链的长度为 $n$，并且可以证明，任何其他极大素理想链的长度也为 $n$. 因而，仅通过对 $\mathcal O(\mathbb C^n)$ 的环论考察，就能恢复 $\mathbb C^n$ 的维数 $n$.
\end{remark*}

\subsection{理想的运算}\label{sec:galois-1-2-9}
设 $I,J$ 是环 $A$ 的两个理想. 我们已经定义了这些理想的和 $I+J$ 与积 $I\cdot J$. 只须回顾，$I+J$ 是由 $I\cup J$ 生成的理想，而 $I\cdot J$ 是由元素 $ij$（$i\in I$，$J\in J$）生成的理想.\footnote{译注：原文写作 $J\in J$，应为 $j\in J$.} 尽管可能引起歧义，我们仍常将 $I\cdot J$ 记为 $IJ$.
\begin{remark*}
若 $I=(a_1,\ldots,a_r)$ 且 $J=(b_1,\ldots,b_s)$，则
\begin{align*}
I+J=(a_1,\ldots,a_r,b_1,\ldots,b_s),
\\
IJ=(a_1b_1,\ldots,a_rb_1,a_1b_2,\ldots,a_rb_2,\ldots,a_rb_s).
\end{align*}
\end{remark*}
当然有理想的包含关系 $I\subset I+J$、$J\subset I+J$ 和 $IJ\subset I\cap J$.
\begin{remark*}
关于理想的包含关系，值得注意的是，它在如下意义下推广了 $A$ 中元素的整除关系：
\begin{align*}
\forall a,b\in A,\quad a\mid b\ \Longleftrightarrow\ (a)\supset(b).
\end{align*}
\end{remark*}
\begin{exercise*}
取 $A=\mathbb Z$、$I=n\mathbb Z$ 和 $J=m\mathbb Z$，并假设 $I,J$ 都是真理想且非零，证明
\begin{align*}
IJ=nm\mathbb Z,\quad I\cap J=\operatorname{lcm}(n,m)\mathbb Z,\quad I+J=\gcd(n,m)\mathbb Z.
\end{align*}
\end{exercise*}
\begin{remark*}
在一般的环 $A$ 中，两个元素 $a,b$ 没有公因子并不意味着 $(a)+(b)=A$. 例如，在 $\mathbb Z[\sqrt{-5}]$ 中，我们已经看到理想 $\mathfrak p=(2)+(1+\sqrt{-5})$ 是真理想，因为 $\mathfrak p^2=(2)$.
\end{remark*}

\subsection{商环}\label{sec:galois-1-2-10}
下面给出一个基本构造，必须充分理解. 设 $A$ 是环，$I$ 是 $A$ 的理想. 在集合 $A$ 上定义等价关系
\begin{align*}
x\equiv y\pmod I\quad\text{当且仅当}\quad x-y\in I.
\end{align*}
这一关系的等价类具有 $x+I=\{x+i,i\in I\}$ 的形式，其中 $x\in A$. 用 $A/I$ 表示等价类的集合，称为 $A$ 关于这一等价关系的商集；用 $\pi_I\colon A\longrightarrow A/I$ 表示把 $x$ 映到其等价类的典范投影. 我们将随意采用 $\pi_I(x)$、$\bar x$、$x+I$ 或 $x\pmod I$ 表示 $x$ 在 $A/I$ 中的像.
\begin{proposition*}
$A/I$ 上存在唯一的交换环结构，使 $\pi_I$ 是环态射.
\end{proposition*}
\begin{proof}
唯一性来自 $\pi_I$ 的满射性. 事实上，要求 $\pi_I$ 是环态射，就迫使有等式
\begin{align*}
\bar x+\bar y=\overline{x+y}\quad\text{和}\quad\bar x\cdot\bar y=\overline{xy}.
\end{align*}
还须说明这一定义是良定义的，并满足环的公理. 为验证良定义性，需要证明：若 $x\equiv x'\pmod I$ 且 $y\equiv y'\pmod I$，则 $(x+y)\equiv(x'+y')\pmod I$ 且 $xy\equiv x'y'\pmod I$. 这是直接的；例如验证第二个关系：写 $x'=x+i$、$y'=y+i'$，其中 $i,i'\in I$，便可看出 $x'y'=xy+i''$，其中 $i''=(iy+i'x+ii')\in I$.
同样，不难验证这样构造的两个运算使 $A/I$ 成为环，其单位元分别为 $\bar0$ 和 $\bar1$.
\end{proof}
称环 $A/I$ 为 $A$ 对 $I$ 的商环.
\begin{example*}
``熟知的''环 $\mathbb Z/n\mathbb Z$ 就是 $\mathbb Z$ 对理想 $(n)=n\mathbb Z$ 的商环.
\end{example*}
\begin{example*}
若 $I=A$，则商环 $A/I$ 是零环.
\end{example*}
$I$ 的性质与 $A/I$ 的性质有如下对应关系.
\begin{proposition*}
设 $A$ 是交换环，$I$ 是 $A$ 的理想.
\begin{enumerate}
\item[i)] $I$ 是极大理想，当且仅当 $A/I$ 是域.
\item[ii)] $I$ 是素理想，当且仅当 $A/I$ 是整环.
\item[iii)] $I$ 是根理想，当且仅当 $A/I$ 是既约环.
\end{enumerate}
\end{proposition*}
\begin{proof}
i) 设 $I$ 极大. 我们要证明 $A/I$ 的每个非零元素都有逆元. 这样的元素具有 $\bar x=x+I$ 的形式，其中 $x\notin I$. 由 $I$ 的极大性，$I+(x)=A$，所以存在 $i\in I$ 和 $y\in A$，使 $i+xy=1$. 于是 $xy\equiv1\pmod I$，所以 $\overline{yx}=\bar1$，从而 $\bar x$ 在 $A/I$ 中确有逆元. 因此 $A/I$ 是域. 反之，设 $A/I$ 是域，$J$ 是严格包含 $I$ 的理想. 须证 $J=A$. 选取元素 $j\in J/I$.\footnote{译注：原文写作 $j\in J/I$；后面的论证实际需要选取 $j\in J\setminus I$.} 它在 $A/I$ 中的像 $\bar j$ 有逆元 $\bar a$，其中 $a\in A$，所以 $aj\in1+I$. 因而 $1\in(j)+I$，于是 $(j)+I=A$，更有 $J=A$.

 ii) 对两个元素 $x,y\in A$，有等价关系 $\bar x=0\Longleftrightarrow x\in I$、$\bar y=0\Longleftrightarrow y\in I$ 和 $\overline{xy}=0\Longleftrightarrow xy\in I$. 因而 ii) 直接由定义得出. iii) 同理.
\end{proof}
\begin{exercise*}
证明映射 $J\mapsto\pi_I^{-1}(J)$ 诱导双射
\begin{align*}
\{A/I\text{ 的理想}\}\xrightarrow{\sim}\{A\text{ 中包含 }I\text{ 的理想}\},
\end{align*}
其逆双射为 $I'\mapsto\pi_I(J)$.\footnote{译注：原文的逆双射写作 $I'\mapsto\pi_I(J)$，变量不一致；应为 $I'\mapsto\pi_I(I')$.} 证明 $\pi_I^{-1}(\sqrt0)=\sqrt I$.
\end{exercise*}


\subsection{商的泛性质}\label{sec:galois-1-2-11}
\begin{proposition*}
对每个满足 $I\subset\operatorname{Ker}(\varphi)$ 的环态射 $\varphi\colon A\longrightarrow A'$，都存在唯一的分解 $\varphi=\bar\varphi\circ\pi_I$，如下图所示.
\[
\begin{tikzcd}
A \arrow[r,"\varphi"] \arrow[d,"\pi_I"'] & A' \\
A/I \arrow[ur,dashed,"\bar\varphi"'] & {}
\end{tikzcd}
\]
此外，$\bar\varphi$ 是单射，当且仅当 $I=\operatorname{Ker}(\varphi)$.
\end{proposition*}
\begin{proof}
唯一性来自 $\pi_I$ 的满射性. 事实上，若 $\bar x\in A/I$，则必有 $\bar\varphi(\bar x)=\varphi(x)$. 为证明存在性，首先须验证这确实无歧义地定义了 $\bar\varphi$. 为此，需要验证若 $x\equiv x'\pmod I$，则 $\varphi(x)=\varphi(x')$. 写 $x'=x+i$，其中 $i\in I$. 于是 $\varphi(x')=\varphi(x)+\varphi(i)=\varphi(x)$，因为 $i\in\operatorname{Ker}(\varphi)$，正如所需. 因而确有映射的分解 $\varphi=\bar\varphi\circ\pi_I$；还须验证 $\bar\varphi$ 是环态射. 但由 $A/I$ 上环结构的定义，这是显然的.

现在证明最后的断言. 先设 $\bar\varphi$ 是单射. 一般事实 $\operatorname{Ker}(\varphi)=\pi_I^{-1}(\operatorname{Ker}(\bar\varphi))$ 表明 $\operatorname{Ker}(\varphi)=\pi_I^{-1}(\{0\})=I$. 反之，设 $\operatorname{Ker}(\varphi)=I$. 对每个 $\bar x\in\operatorname{Ker}(\bar\varphi)$，由同一一般事实可知 $x\in\operatorname{Ker}(\varphi)$，所以 $x\in I$，且 $\bar x=\bar0=0$.
\end{proof}
\begin{remark*}[{什么是泛性质？}]
在``元数学''语言中，若一个对象的某个性质能在同一``范畴''的所有对象中刻画这个对象，就称这一性质为泛性质. 换言之，具有这一泛性质的对象在唯一同构的意义下唯一确定（这里的同构指所考察对象范畴中适用的同构）. 在这里，二元组 $(A/I,\pi)$ 属于满足 $\psi(I)=0$ 的 $A$-代数 $(B,\psi)$ 的``范畴''，上述命题说明它具有如下性质：对每个满足 $\varphi(I)=0$ 的 $A$-代数 $(A',\varphi)$，存在唯一的 $A$-代数态射 $A/I\longrightarrow A'$. 假设另一个满足 $\psi(I)=0$ 的 $A$-代数 $(B,\psi)$ 具有同一性质. 对 $(A',\varphi)=(A/I,\pi_I)$ 应用这一性质，得到一个 $A$-代数态射 $B\xrightarrow{\bar\pi}A/I$. 另一方面，上述命题给出一个 $A$-代数态射 $A/I\xrightarrow{\bar\psi}B$. 复合 $\bar\pi\bar\psi$ 是 $A$-代数 $A/I$ 的自同态，而上述命题告诉我们，这样的自同态只有一个. 恒等映射显然是 $A$-代数自同态，所以必有 $\bar\pi\circ\bar\psi=\operatorname{id}_{A/I}$. 同样，由所假设的 $(B,\psi)$ 的性质可知 $\bar\psi\circ\bar\pi=\operatorname{id}_B$，因此得到一个 $A$-代数同构 $A/I\xrightarrow{\sim}B$；由所假设的性质或上述命题，这一同构还是唯一的.
\end{remark*}
\begin{corollary*}
每个环态射 $\varphi\colon A\longrightarrow A'$ 都有唯一的分解
\begin{align*}
\varphi\colon A\twoheadrightarrow A/\operatorname{Ker}(\varphi)
\xrightarrow{\sim}\operatorname{Im}(\varphi)\hookrightarrow A',
\end{align*}
其中第一个箭头是到商 $A/\operatorname{Ker}(\varphi)$ 的典范投影.
\end{corollary*}
回顾，符号 $\twoheadrightarrow$ 表示满射，符号 $\hookrightarrow$ 表示单射，符号 $\xrightarrow{\sim}$ 表示同构.
\begin{proof}
对 $I=\operatorname{Ker}(\varphi)$ 应用上述命题，得到分解 $\varphi=\bar\varphi\circ\pi_I$. 此外，$\bar\varphi$ 是单射，所以它给出其源到其像 $\operatorname{Im}(\bar\varphi)=\operatorname{Im}(\varphi)$ 的同构，而这个像是 $A'$ 的一个子环.
\end{proof}
\begin{example*}
设 $V$ 是 $\mathbb C^n$ 的一个代数子集. 由多项式函数代数 $\mathcal O(V)$ 的定义，函数的限制映射 $\mathcal O(\mathbb C^n)\longrightarrow\mathcal O(V)$ 是满射. 设 $\mathcal I$ 是它的理想，\footnote{译注：原文仅说``它的理想''（son idéal），未明确指出是核理想；上下文所需的正是限制映射的核.} 因而它诱导同构 $\mathcal O(\mathbb C^n)/\mathcal I\xrightarrow{\sim}\mathcal O(V)$，将 $\mathcal O(V)$ 表示为多项式代数 $\mathbb C[X_1,\ldots,X_n]$ 的一个商.
\end{example*}

\subsection{商之间的态射}\label{sec:galois-1-2-12}
现在设 $J\supset I$ 是 $A$ 中包含 $I$ 的理想. 商 $A/I$ 的泛性质给出分解
\begin{align*}
\pi_J\colon A\xrightarrow{\pi_I}A/I\xrightarrow{\bar\pi_J}A/J.
\end{align*}
\begin{proposition*}
$J$ 在 $A/I$ 中的像 $J/I:=\pi_I(J)$ 是理想，且态射 $\bar\pi_J$ 诱导同构
\begin{align*}
(A/I)/(J/I)\xrightarrow{\sim}A/J.
\end{align*}
\end{proposition*}
\begin{proof}
由于态射 $\pi_J$ 是满射，态射 $\bar\pi_J$ 也是满射，所以只须说明其核为 $\operatorname{Ker}(\bar\pi_J)=\pi_I(J)$（这也会同时证明 $\pi_I(J)$ 确实是理想）. 有
\begin{align*}
\pi_I^{-1}(\operatorname{Ker}(\bar\pi_J))
=\operatorname{Ker}(\bar\pi_J\circ\pi_I)
=\operatorname{Ker}(\pi_J)=J.
\end{align*}
但由于 $\pi_I$ 是满射，故
\footnote{译注：原文此式在中间表达式末尾少一个右圆括号，照录其括号形式.}
\begin{align*}
\operatorname{Ker}(\bar\pi_J)
=\pi_I(\pi_I^{-1}(\operatorname{Ker}(\bar\pi_J))
=\pi_I(J).
\end{align*}
\end{proof}
\begin{example*}
由此重新得到如下``熟知的''事实：当 $m\mid n$ 时，映射 $a\mapsto a\pmod m$ 经由映射 $a\mapsto a\pmod n$ 分解为一个态射 $\mathbb Z/n\mathbb Z\longrightarrow\mathbb Z/m\mathbb Z$，并诱导同构
\begin{align*}
(\mathbb Z/n\mathbb Z)/m(\mathbb Z/n\mathbb Z)\xrightarrow{\sim}\mathbb Z/m\mathbb Z.
\end{align*}
\end{example*}
一个变式：不从 $J\supset I$ 出发，而取任意 $J$，并对包含 $I$ 的理想 $I+J$ 应用上述命题. 得到分解 $\pi_{I+J}=\bar\pi_{I+J}\circ\pi_I$，其中 $\bar\pi_{I+J}$ 诱导同构
\begin{align*}
(A/I)/((I+J)/I)\xrightarrow{\sim}A/(I+J).
\end{align*}
\subsection{中国剩余定理}\label{sec:galois-1-2-13}
设 $I,J$ 是 $A$ 的两个理想. 态射
\begin{align*}
A\longrightarrow A/I\times A/J,\qquad a\longmapsto(a\pmod I,a\pmod J)
\end{align*}
的核显然包含 $I\cap J$，所以由泛性质，该态射经由 $\pi_{I\cap J}$ 分解为态射
\begin{align*}
A/(I\cap J)\xrightarrow{(\bar\pi_I,\bar\pi_J)}A/I\times A/J.
\end{align*}
\begin{proposition*}
上述态射 $(\bar\pi_I,\bar\pi_J)$ 是单射. 它是满射，当且仅当 $I+J=A$. 因而在后一情形下得到同构
\begin{align*}
A/(I\cap J)\xrightarrow{\sim}A/I\times A/J.
\end{align*}
\end{proposition*}
\begin{proof}
为证明单射性，只须验证 $I\cap J$ 恰为态射 $(\pi_I,\pi_J)$ 的核，这是显然的.

为证明 $(\bar\pi_I,\bar\pi_J)$ 的满射性，需要说明对所有 $a,b\in A$，交集 $(a+I)\cap(b+J)$ 非空. 事实上，若如此，则对这一交集中的任意 $c$ 都有
\begin{align*}
(c\pmod I,c\pmod J)=(a\pmod I,b\pmod J).
\end{align*}
因此，设 $I+J=A$，选取 $i\in I$、$j\in J$，使 $i+j=1$. 那么元素 $c=aj+bi$ 属于 $a+I$，因为 $c=a-ai+bi$；它也属于 $b+J$，因为 $c=b-bj+aj$. 因而 $(\bar\pi_I,\bar\pi_J)$ 确实是满射.

反之，设 $(\bar\pi_I,\bar\pi_J)$ 是满射. 则特别地，存在 $j\in A$，使 $(\bar1,0)=(\pi_I(j),\pi_J(j))$；这等价于 $j\in1+I$ 且 $j\in J$. 令 $i=1-j$，则 $i\in I$ 且 $i+j=1$，故得到所需的 $I+J=A$.
\end{proof}
\begin{example*}
取 $A=\mathbb Z$、$I=n\mathbb Z$、$J=m\mathbb Z$，并假设 $I,J$ 都是真理想且非零，则有 $I\cap J=\operatorname{lcm}(n,m)\mathbb Z$ 和 $I+J=\gcd(n,m)\mathbb Z$. 特别地，条件 $I+J=A$ 等价于 $\gcd(n,m)=1$；此时重新得到通常的中国剩余引理：$\mathbb Z/nm\mathbb Z\xrightarrow{\sim}\mathbb Z/n\mathbb Z\times\mathbb Z/m\mathbb Z$.
\end{example*}

